Según las fuentes periodísticas citadas, el jueves 8 de octubre de 2026, cerca de las 17:00 (UTC$-$3), un flujo de agua, barro y rocas descendió desde la quebrada que drena el sector precordillerano de San Carlos de Apoquindo (Las Condes, Santiago) y recorrió calles residenciales por varios kilómetros. Arrastró vehículos y un bus \citep{theclinic2026frontal}, inundó viviendas, una clínica y una universidad, dejó seis heridos y motivó la declaración de zona de catástrofe \citep{interferencia2026}. Según la información periodística consultada, que recoge antecedentes del Servicio Nacional de Geología y Minería (Sernageomin), el evento ocurrió en una zona clasificada en febrero de 2026 como de alto peligro por flujos de detritos \citep{chilevision2026}. Según el municipio y BioBioChile, la cuenca había sido afectada por un incendio de unas \SI{1000}{ha} entre diciembre de 2025 y enero de 2026 \citep{lascondes2026,biobio2026origen}. Según el CR2, el evento siguió a un comienzo de octubre excepcionalmente lluvioso: más de \SI{90}{mm} en Quinta Normal en lo que iba del mes, cerca del máximo histórico de octubre de 1930 \citep{cr2_2026}.

Al día siguiente, entre las fuentes públicas consultadas no se había localizado una cartografía oficial del área afectada, ni aforos, ni marcas de nivel publicadas. La información disponible se reducía a notas de prensa, la lista de calles cerradas por el municipio, un perímetro de afectación y dos versiones de las fuentes sobre el origen que no coinciden entre sí. Este trabajo ensaya escenarios de recorrido, volumen y caudal del aluvión mediante un modelo físico explícito y evidencia pública verificable. Asimismo, calcula para esos escenarios la profundidad, la velocidad y la presión del flujo en calles y edificios del área de estudio, considerando la incertidumbre de las estimaciones.
